\chapter{Urinary Urgency}

\section*{Definition}
Urinary urgency is a sudden, difficult-to-defer need to urinate. It may occur with urinary frequency, nighttime urination, leakage, or pain. Urgency is a symptom rather than a diagnosis and may arise from infection, bladder irritation, overactive bladder, obstruction, or neurologic disease.

\section*{What Happens in Your Body}
As the bladder fills, coordinated nerve signals normally allow storage until urination is convenient. In urgency, involuntary bladder contractions, inflammation, increased sensory signaling, reduced bladder capacity, or impaired control create an urgent sensation. Neurologic disease can disrupt communication between the bladder, spinal cord, and brain.

\section*{Causes}
\begin{itemize}
  \item Urinary tract infection, urethritis, vaginitis, or bladder irritation
  \item Overactive bladder, pelvic-floor dysfunction, or bladder pain syndrome
  \item Enlarged prostate, urethral narrowing, stone, or other obstruction
  \item Diabetes, pregnancy, constipation, stroke, Parkinson disease, or multiple sclerosis
  \item Caffeine, alcohol, excessive fluids, diuretics, and some medications
\end{itemize}

\section*{When to See a Doctor}
Seek prompt care for urgency with fever, flank pain, vomiting, visible blood, pregnancy, severe pain, or inability to urinate. Emergency assessment is needed for new bladder symptoms with leg weakness, saddle numbness, or loss of bowel control.

\section*{Related Symptoms}
\begin{itemize}
  \item Frequency, nocturia, urge leakage, hesitancy, weak stream, or incomplete emptying
  \item Burning urination, lower abdominal pain, cloudy or foul-smelling urine
\end{itemize}

\section*{Self-Care / Home Management}
Drink a normal amount of water unless a clinician has prescribed fluid restriction. Reduce caffeine, alcohol, carbonated drinks, and known bladder irritants. Keep a bladder diary and avoid prolonged holding, but do not self-treat suspected infection with leftover antibiotics.

\section*{How Your Doctor Will Evaluate You}
\begin{itemize}
  \item History covers timing, leakage, pain, fever, sexual and pregnancy status, bowel function, neurologic symptoms, and medications.
  \item Urinalysis and urine culture are used when infection is possible; pregnancy testing and glucose testing may be appropriate.
  \item Examination may include abdomen, pelvis, prostate, neurologic function, and post-void residual volume.
  \item Ultrasound, cystoscopy, or urodynamic testing is reserved for recurrent, complicated, or unexplained symptoms.
\end{itemize}

\section*{Treatment Options}
Treatment may include antibiotics for confirmed infection, bladder training, pelvic-floor therapy, constipation treatment, medication for overactive bladder, or management of obstruction and neurologic disease. Persistent hematuria or recurrent symptoms require evaluation for structural disease.

\section*{References}
\begin{thebibliography}{99}
\bibitem{urinary_urgency:ref1} Gormley EA, et al. Diagnosis and treatment of overactive bladder. \textit{J Urol}. 2019;202:558--563.
\bibitem{urinary_urgency:ref2} Lightner DJ, et al. AUA/SUFU Guideline on the Diagnosis and Treatment of Idiopathic Overactive Bladder. AUA; 2024.
\bibitem{urinary_urgency:ref3} Hooton TM, et al. Diagnosis and treatment of uncomplicated urinary tract infection. \textit{Clin Infect Dis}. 2010;52:e103--e120.
\end{thebibliography}
